% Standalone appendix document - compiles on its own:
%   latexmk -pdf -outdir=build 08-agenda-2026-11-04.tex
\documentclass[a4paper,11pt]{article}
\usepackage[danish]{babel}
\usepackage{../../appendix-style}

\title{12.8 Dagsorden -- 04-11}
\author{SW4PRJ4 -- Group 3}
\date{04-11-2026}

\begin{document}
\maketitle

\section*{Mødeinformation}

\textbf{Dato:} 04-11-2026 (onsdag)

\textbf{Tid:} 10:00

\textbf{Sted:} Incuba

\textbf{Deltagere:} Gruppe 3 (internt møde)

\section*{Dagsordenspunkter}

\begin{enumerate}
    \item \textbf{Opfølgning på vejledermødet 3/11}
    \begin{itemize}
        \item Hvad svarede Michel om diagrammer, om kode eller rapport, om sprint 5 og om at droppe should-haves? Se referatet fra 3/11.
    \end{itemize}
    \item \textbf{Sprint review for sprint 4 (28/10 til 3/11)} (ca. 30 min)
    \begin{itemize}
        \item Sprint goal fra 28/10 var at gøre kortene fra sprint 3 færdige og få rapporten op på niveau med koden med diagrammer og gennemskrevne dokumenter. Nåede vi det?
        \item Vi viser login og onboarding sammen med diagrammerne. Demoen blev udskudt fra 28/10.
        \item Bagefter gennemgår vi login med Keycloak ud fra diagrammet, så alle forstår, hvordan det hænger sammen med resten af systemet.
        \item 16 kort fra sprint 4 venter på godkendelse i \emph{In Review}, i alt 37 point. Efter demoen godkender vi dem og flytter dem til \emph{Done}.
        \item Tre kort er ikke færdige. Databasen (5 point) og integrationsplanen (2 point) er i gang. Netværksarkitekturen (1 point) er ikke startet.
        \item Integrationsplanen og et dependency tree blev merget 28/10 (PR \#40). Hvad mangler der på kortet?
        \item Kortet med sprint goals og velocity (1 point) er færdigt, når tallene for sprint 4 er skrevet ind.
        \item Kortet om at huske login på telefonen (5 point) venter på godkendelse, men har ingen sprint-label. Tæller den med i sprint 4?
        \item Kortet om sproget i apps'ene har ingen beskrivelse, point eller labels. Vi har allerede valgt engelsk, så skal vi lukke den?
        \item Vi committede 46,5 point i sprint 4, og hele backloggen er nu på 186 point. Godkender vi alt i \emph{In Review}, og bliver kortet med velocity færdigt, er vi færdige med 38,5 point, og 8 point går videre. Tallene skrives i velocity-filen.
    \end{itemize}
    \item \textbf{Retrospective} (ca. 20 min)
    \begin{itemize}
        \item Sidst aftalte vi at bruge dependency trees og det fra undervisningen i softwaretest, når vi skriver tests. Gjorde vi det?
        \item Vi aftalte også, at hver person selv flytter sine kort på boardet. Gjorde vi det?
        \item Databasen var lovet til vejledermødet. Nåede vi det, og hvis ikke, hvorfor?
        \item Hvad gik godt?
        \item Hvad gik skidt?
        \item Hvad ændrer vi i sprint 5? Hver ting bliver et kort i Kaneo.
    \end{itemize}
    \item \textbf{Sprint planning for sprint 5 (4/11 til 17/11)} (ca. 45 min)
    \begin{itemize}
        \item Sprint 5 er den sidste sprint med nye funktioner, og vi planlægger kun must-have-krav. Skal planen rettes efter det, Michel sagde?
        \item Hvor mange timer har hver af os de næste to uger? I sprint 1 til 3 blev vi i gennemsnit færdige med 27,8 point.
        \item Kort, der går videre fra sprint 4, beholder label \emph{Sprint 4} og får \emph{Sprint 5} oveni.
        \item Kortet om zombie-brugere er stadig ikke estimeret. Estimerer vi den, eller dropper vi den?
        \item Kort fra backloggen, vi kan tage ind:
        \begin{itemize}
            \item Standard-pakketyper (5 point) og udløbsdato og næste ordre (5 point).
            \item Tilmelding til abonnement (3 point) og kundens første ordre (5 point).
            \item Listen over ordrer i admin-panelet (3 point) og at markere en ordre som sendt (2 point).
            \item C4 Context-diagram (1 point), deploymentdiagram (2 point), kapitel 3 (2 point) og kapitel 4 (3 point).
        \end{itemize}
        \item Et kort må kun flyttes til \emph{To Do}, når det er klar. Det kræver kravnummeret som første linje, definition of done, point, epic og en ejer.
        \item Vi skriver et sprint goal på én sætning, der nævner de krav, sprinten leverer.
        \item Hvert kort får en ejer og en undergruppe.
        \item Vi lægger point i \emph{To Do} sammen, så vi har committed.
    \end{itemize}
    \item \textbf{Eventuelt}
\end{enumerate}

\end{document}
